\chapter{绪论}

\section{研究背景}

这里填写论文的研究背景。正文段落默认首行缩进 2 字符，字号为小四，行距为固定 20 pt。使用模板时，可直接将本文件替换为自己的第一章内容。

文献引用可以使用普通引用命令 \cite{hartley2004multiple}，也可以使用上标引用命令 \upcite{lepetit2009epnp}。参考文献由 \texttt{refs.bib} 和 \texttt{gbt7714-numerical} 样式统一生成。

\section{研究意义}

这里填写研究意义、研究目标或应用价值。建议每一节围绕一个明确主题展开，避免标题层级过深。

\section{本文结构}

本文结构可按如下方式组织：第一章介绍研究背景与意义；第二章给出相关理论、方法或系统设计；第三章总结全文工作并展望后续研究。
